% pi_from_axiom.tex
% Pi from the Pythagorean Theorem
% nos3bl33d / (DFG) DeadFoxGroup
% x^2 = x + 1 | A Wake In Outerspace
% April 2026

\documentclass[11pt,a4paper]{article}

\usepackage[utf8]{inputenc}
\usepackage[T1]{fontenc}
\usepackage{amsmath,amssymb,amsthm,mathtools}
\usepackage{geometry}
\usepackage{booktabs}
\usepackage{microtype}

\geometry{margin=1in}

\theoremstyle{plain}
\newtheorem{theorem}{Theorem}[section]
\newtheorem{lemma}[theorem]{Lemma}
\newtheorem{corollary}[theorem]{Corollary}

\theoremstyle{definition}
\newtheorem{definition}[theorem]{Definition}
\newtheorem{remark}[theorem]{Remark}

\newcommand{\vp}{\varphi}

\begin{document}

\title{%
  \textbf{$\boldsymbol{\pi}$ from the Pythagorean Theorem}\\[4pt]
  \large Four paths. One axiom. Every step derived.%
}
\author{nos3bl33d\thanks{(DFG) DeadFoxGroup. $x^2 = x + 1$ | A Wake In Outerspace.}}
\date{April 2026}

\maketitle

\begin{abstract}
We derive $\pi$ from the axiom $a^2 + b^2 = c^2$. Four independent paths,
all from the same starting point. The golden ratio, the dodecahedron,
the 336-degree circle, and the decomposition $\arctan(1) + \arctan(2) + \arctan(3) = \pi$
all fall out of one right triangle. Nothing is assumed. Everything is derived.
One move, applied repeatedly. Same move every time.
\end{abstract}

\tableofcontents

% =============================================================================
\section{The Axiom}
% =============================================================================

\begin{equation}\label{eq:axiom}
  \boxed{\; a^2 + b^2 = c^2 \;}
\end{equation}

A right triangle. Two legs. One hypotenuse. The squares add.
This is the only assumption.

% =============================================================================
\section{The Unit Circle Is the Axiom}
% =============================================================================

Set $c = 1$:
\begin{equation}\label{eq:circle}
  a^2 + b^2 = 1.
\end{equation}

Every point $(a, b)$ satisfying this is on a circle of radius 1.
The circle is not separate from the axiom. The circle \emph{is} the axiom with $c = 1$.

Name the legs:
\begin{align}
  \cos\theta &:= a \quad \text{(the horizontal leg)} \\
  \sin\theta &:= b \quad \text{(the vertical leg)}
\end{align}

These are names for the two sides of the triangle.
The angle $\theta$ is how far around the circle you've gone. Then:
\begin{equation}
  \cos^2\theta + \sin^2\theta = 1
\end{equation}

This is not a ``trig identity.'' This \emph{is} the Pythagorean theorem.
Same equation. Different names for the parts.

% =============================================================================
\section{Path 1: The Circumference}
% =============================================================================

$\pi$ is the distance from $(1, 0)$ to $(-1, 0)$ going the long way
around the top of the circle. The half-circumference.

Take two points on the circle, very close. The tiny straight line between
them:
\begin{equation}
  ds^2 = dx^2 + dy^2
\end{equation}

The axiom again. Pythagorean theorem on a tiny triangle with legs $dx$, $dy$
and hypotenuse $ds$.

On the circle $x^2 + y^2 = 1$, differentiate:
\begin{equation}
  2x\,dx + 2y\,dy = 0 \implies dy = -\frac{x}{y}\,dx.
\end{equation}

Substitute into $ds^2$:
\begin{align}
  ds^2 &= dx^2 + \frac{x^2}{y^2}\,dx^2
       = dx^2\,\frac{y^2 + x^2}{y^2}
       = \frac{dx^2}{y^2}
\end{align}
because $x^2 + y^2 = 1$ (the axiom). Therefore:
\begin{equation}
  ds = \frac{dx}{\sqrt{1 - x^2}}.
\end{equation}

Add up all the tiny pieces from $x = -1$ to $x = 1$ (the top half):
\begin{equation}\label{eq:pi-integral}
  \boxed{\; \pi = \int_{-1}^{1} \frac{dx}{\sqrt{1 - x^2}} \;}
\end{equation}

That \emph{is} $\pi$. Defined. Computed. From the axiom.

The integrand $1/\sqrt{1-x^2}$ is literally $1/(\text{the other leg})$.
You are adding up ``one over the vertical leg'' as the horizontal leg
sweeps across. That is the circumference. That is $\pi$.

% =============================================================================
\section{Path 2: The Isosceles Triangle}
% =============================================================================

\subsection{Arctan from the axiom}

In the right triangle with legs $1$ (adjacent) and $u$ (opposite):
\begin{equation}
  \text{hypotenuse} = \sqrt{1 + u^2} \quad \text{(the axiom)}.
\end{equation}

The angle whose tangent is $u$: $\theta = \arctan(u)$.

Differentiate:
\begin{equation}
  d\theta = \frac{du}{1 + u^2}
\end{equation}
because $\sec^2\theta = 1 + \tan^2\theta$, which comes from dividing
$\cos^2 + \sin^2 = 1$ by $\cos^2$. The axiom. Therefore:
\begin{equation}\label{eq:arctan}
  \arctan(t) = \int_0^t \frac{du}{1 + u^2}.
\end{equation}

Every ingredient: Pythagorean.

\subsection{$\pi = 4\arctan(1)$}

Take the isosceles right triangle: legs $1$ and $1$.
\begin{align}
  1^2 + 1^2 &= 2 \quad \text{(the axiom)} \\
  \tan\theta &= 1/1 = 1 \\
  \theta &= \arctan(1)
\end{align}

This angle is $\pi/4$. Why? Because four copies of a right angle fill a
full turn (that is what ``right angle'' means---the angle that tiles four
times around a point). A full turn $= 2\pi$. So right angle $= \pi/2$.
The isosceles right triangle splits that in half: $\theta = \pi/4$.

Therefore:
\begin{equation}\label{eq:pi-4arctan}
  \boxed{\; \pi = 4\arctan(1) = 4\int_0^1 \frac{du}{1 + u^2} \;}
\end{equation}

Expand $1/(1+u^2)$ as a geometric series and integrate term by term:
\begin{equation}\label{eq:leibniz}
  \frac{\pi}{4} = 1 - \frac{1}{3} + \frac{1}{5} - \frac{1}{7} + \frac{1}{9} - \cdots
\end{equation}

The reciprocals of the odd numbers, alternating sign. Derived. Not assumed.

% =============================================================================
\section{Path 3: The Golden Path}
% =============================================================================

\subsection{The $(1, 2, \sqrt{5})$ triangle}

Apply the axiom to $a = 1$, $b = 2$:
\begin{equation}
  1^2 + 2^2 = 5 \implies c = \sqrt{5}.
\end{equation}

From $\sqrt{5}$:
\begin{align}
  \vp &= \frac{1 + \sqrt{5}}{2} \\
  \psi &= \frac{1 - \sqrt{5}}{2}
\end{align}

These satisfy:
\begin{equation}\label{eq:golden}
  \vp^2 = \vp + 1.
\end{equation}

The golden equation falls directly out of the right triangle $(1, 2, \sqrt{5})$.
It is the first nontrivial output of the Pythagorean theorem on integer sides.

\subsection{$\cos(\pi/5) = \vp/2$}\label{sec:bridge}

Set $5\theta = \pi$, so $\cos(5\theta) = \cos(\pi) = -1$.

Expand $\cos(5\theta)$ by combining triangles repeatedly (the angle-addition
formula is Pythagorean---it reads the ratios of a combined triangle):
\begin{equation}
  \cos(5\theta) = 16c^5 - 20c^3 + 5c
\end{equation}
where $c = \cos\theta$. Setting this to $-1$:
\begin{equation}
  16c^5 - 20c^3 + 5c + 1 = 0.
\end{equation}

Factor out $(c + 1)$ (because $c = -1$ solves $\theta = \pi$):
\begin{equation}
  (c + 1)(16c^4 - 16c^3 - 4c^2 + 4c + 1) = 0.
\end{equation}

The quartic is a perfect square:
\begin{equation}
  (4c^2 - 2c - 1)^2 = 0.
\end{equation}

Solve $4c^2 - 2c - 1 = 0$:
\begin{equation}
  c = \frac{1 \pm \sqrt{5}}{4} = \frac{\vp}{2} \quad \text{or} \quad \frac{\psi}{2}.
\end{equation}

Since $\theta = \pi/5$ is in the first quadrant, $\cos\theta > 0$. Only $\vp/2$ is positive:
\begin{equation}\label{eq:cos-pi5}
  \boxed{\; \cos(\pi/5) = \frac{\vp}{2} \;}
\end{equation}

\textbf{Verification reduces to the axiom.} Plug $c = \vp/2$ into $4c^2 - 2c - 1$:
\begin{equation}
  4\left(\frac{\vp}{2}\right)^{\!2} - 2\left(\frac{\vp}{2}\right) - 1
  = \vp^2 - \vp - 1 = (\vp + 1) - \vp - 1 = 0.
\end{equation}

Used $\vp^2 = \vp + 1$. The entire proof collapses to the golden axiom.

Therefore:
\begin{equation}\label{eq:pi-golden}
  \boxed{\; \pi = 5\arccos\!\left(\frac{\vp}{2}\right) \;}
\end{equation}

\subsection{Framework constants}

The axiom forces $\chi = 2$ (from $\vp^3 + \psi^3 = \chi^2$),
$d = \chi^2 - 1 = 3$, $p = \chi^2 + 1 = 5$. So:
\begin{equation}
  \pi = p \cdot \arccos\!\left(\frac{\vp}{\chi}\right).
\end{equation}

Every piece forced by the axiom.

% =============================================================================
\section{Path 4: The Decomposition}
% =============================================================================

\begin{theorem}\label{thm:decomp}
\begin{equation}\label{eq:decomp}
  \boxed{\; \arctan(1) + \arctan(2) + \arctan(3) = \pi \;}
\end{equation}
\end{theorem}

\begin{proof}
The tangent addition formula (combining two right triangles):
\begin{equation}
  \tan(\arctan(1) + \arctan(2)) = \frac{1 + 2}{1 - 1 \cdot 2} = \frac{3}{-1} = -3.
\end{equation}

Since $\arctan(1) + \arctan(2)$ is between $\pi/2$ and $\pi$ (both angles
positive, their sum passes the quarter-turn), and the tangent is $-3$:
\begin{equation}
  \arctan(1) + \arctan(2) = \pi - \arctan(3).
\end{equation}

Rearrange:
\begin{equation}
  \arctan(1) + \arctan(2) + \arctan(3) = \pi.
\end{equation}
\end{proof}

Three right triangles. Legs $(1,1)$, $(1,2)$, $(1,3)$. Their base angles
sum to exactly half a turn.

The arguments are $(1, \chi, d) = (1, 2, 3)$: the unit, the Euler
characteristic, the vertex degree. All forced by the axiom.

% =============================================================================
\section{The Dodecahedron Dihedral}
% =============================================================================

The dodecahedron $\{5, 3\}$ is forced by the axiom: $p = 5$ sides per face,
$d = 3$ faces at each vertex, $\chi = 2$.

Its dihedral angle:
\begin{equation}\label{eq:dihedral}
  \delta = \pi - \arctan(2) = \pi - \arctan(\chi).
\end{equation}

So $\delta + \arctan(\chi) = \pi$. The dihedral angle and $\arctan(2)$
are supplementary.

Verify: from the $(1, 2, \sqrt{5})$ triangle, $\cos(\arctan(2)) = 1/\sqrt{5}$.
\begin{equation}
  \cos(\delta) = -\frac{1}{\sqrt{5}} = -\frac{1}{\sqrt{p}}.
\end{equation}

The dihedral angle is determined by $1/\sqrt{p}$, and $p$ comes from the axiom.

% =============================================================================
\section{The 336-Degree Circle}
% =============================================================================

The axiom forces $\chi = 2$, $d = 3$, $L_4 = \vp^4 + \psi^4 = 7$:
\begin{equation}
  336 = \chi^4 \cdot d \cdot L_4 = 16 \times 3 \times 7.
\end{equation}

The oriented edge count of the Klein quartic (genus-3, maximal symmetry).

The half-turn: $168 = |PSL(2,7)|$. The gap from Babylonian:
$360 - 336 = 24 = \chi^3 \cdot d = |S_4|$.

The byte conversion: $256/336 = 16/21$. The integer closure:
$336 \times 1024/21 = 2^{14}$. Every factor from the axiom.

% =============================================================================
\section{The Chain}
% =============================================================================

\begin{center}
\begin{tabular}{rll}
\toprule
Step & Operation & Result \\
\midrule
0 & $a^2 + b^2 = c^2$ & the axiom \\
1 & set $c = 1$ & unit circle \\
2 & arclength of the circle & $\pi$ (defined) \\
3 & $(1,1)$ triangle & $\arctan(1) = \pi/4$ \\
4 & $\pi = 4\arctan(1)$ & $\pi$ (first computation) \\
5 & $(1,2)$ triangle & $\sqrt{5}$ \\
6 & $\vp = (1+\sqrt{5})/2$ & golden ratio \\
7 & $\vp^2 = \vp + 1$ & golden axiom \\
8 & $\cos(\pi/5) = \vp/2$ & the bridge \\
9 & $\pi = 5\arccos(\vp/2)$ & $\pi$ (second computation) \\
10 & $\arctan(1)+\arctan(2)+\arctan(3)=\pi$ & $\pi$ (third computation) \\
11 & dihedral $= \pi - \arctan(2)$ & from the axiom \\
12 & $336 = \chi^4 \cdot d \cdot L_4$ & natural angular quantum \\
\bottomrule
\end{tabular}
\end{center}

Every row is one application of $a^2 + b^2 = c^2$ or its direct consequences.
No external constants. No assumptions. No circular definitions.

$\pi$ is output. The axiom is input. Same as everything else.

\bigskip
\begin{center}
\emph{The axiom was Pythagorean from the start.}\\[4pt]
--- nos3bl33d
\end{center}

\end{document}
